%% The abstract environment is used for the abstract of a chapter.
\newenvironment{chapterabstract}{%
    \list{}{\leftmargin\rightmargin}%
    \item%
    \relax%
    \rmfamily\itshape%
}{%
    \endlist%
}

\chapter{Introduction} \pagestyle{intro-biblio}
\chapterprecishere{
    \begin{epigraphs}
        \qitem{
            \textit{Nature and nature's laws lay hid in the night; God said `Let Newton be!' and all was light.}
        }{\opensans Alexander Pope} \\[1\baselineskip]
        \qitem{
            \textit{It did not last: the devil shouting `Ho. Let Einstein be!' restore the status quo.}
        }{\opensans Sir John Collings Squire}
    \end{epigraphs}
    
    \vspace{0pt plus 1filll} % fill space up to bottom

    \begin{chapterabstract}
        \lipsum[1]
    \end{chapterabstract}

    \blfootnote{Parts of this chapter have been published in Annalen der Physik \textbf{324}, 289 (1906)}
}
\chapterprecistoc{This command allows you to add a brief abstract to the ToC, in case you are interested in that sort of stuff. If not, then either leave this field empty, or don't bother calling upon this command.}

\lipsum[0-12]

